% !TEX program = xelatex
\documentclass[11pt,a4paper]{ctexart}

\usepackage[margin=2.4cm]{geometry}
\usepackage{amsmath,amssymb}
\usepackage{booktabs}
\usepackage{longtable}
\usepackage{array}
\usepackage{enumitem}
\usepackage{hyperref}
\usepackage{xcolor}
\usepackage{caption}
\hypersetup{colorlinks=true,linkcolor=black,urlcolor=blue,citecolor=black}
\captionsetup{font=small,labelfont=bf}
\setlength{\parskip}{0.4em}
\setlength{\emergencystretch}{3em}
\setlist{topsep=0.3em,itemsep=0.2em,parsep=0.15em}

\title{\bfseries 《AI 在金融交易中的应用：一份可落地的研究综述》\\ 同行评审报告\\\large 重点审理：复现协议与证据表的严谨性}
\author{\large 评审人：叨客宁宁}
\date{}

\begin{document}
\maketitle

\section*{评审基本信息与配置}
\addcontentsline{toc}{section}{评审基本信息与配置}

\begin{table}[htbp]
\centering
\begin{tabular}{@{}>{\raggedright\arraybackslash}p{3.4cm}>{\raggedright\arraybackslash}p{11.0cm}@{}}
\toprule
被评审稿件 & 《AI 在金融交易中的应用：一份可落地的研究综述》（61 页，21 章） \\
评审模式 & 多视角独立评审 + 系统清单 + 编辑综合 \\
审稿视角 & 主编、方法学、领域、跨学科与实际影响、Devil's Advocate（共 5 个） \\
评审范围 & 重点：第 12 章证据表、第 13 章比较矩阵、第 15 章复现核验、第 16 章可复现协议；兼及全篇 \\
约束声明 & 只读评审，不修改稿件；每条批评均给出“错在哪、在何处、怎么改” \\
\bottomrule
\end{tabular}
\end{table}

\noindent\textbf{评审纪律说明}：五个视角独立评审，互不参照；编辑综合阶段的每一条结论均溯源到具体视角报告，未引入任何视角之外的新意见。Devil's Advocate 的意见单独列示，不与其他视角混合。

\section{阶段一：五视角独立评审}

\subsection{视角 A · 主编（期刊契合与整体质量）}

\textbf{中心问题与主要发现}。稿件试图回答“AI 在金融交易中到底改变了什么”，并给出一条贯穿全篇的判据——可交易缝隙 $G(\pi)=\mathbb{E}[r_\pi]-C(\pi)-\lambda_\pi$。这个选择的优点是明显的：它把“方法好不好”这个伪问题，替换成“扣掉摩擦与选择偏差后还剩多少”这个可操作的问题。第 2 章给出的缝隙存在必要条件（$\delta_0>c$ 当且仅当存在正收益调仓周期）是全篇最漂亮的段落：推导只有两行，却同时解释了高频策略对成本的敏感性、以及低频信号被迫降频的现象。第 11 章的成本分解（日频 $5$bp 毛收益对 $10$bp 成本）把抽象判据落成了可算的数，写作上是有意识的“先立判据、再给实例”。

\textbf{主要问题}。稿件的定位声明与呈现方式之间存在张力。第 1 章把本文与既有综述的差异归纳为“把评测本身当作研究对象”“区分架构事实与判断结论”，但作为核心证据载体的第 12 章证据表，其列结构与这一声明不匹配——证据表只有 ID、标题、来源、年、关键定位、核验状态六列，既不包含样本外检验，也不包含复现状态，而这两者恰恰是“评测严谨性”的直接体现（详见方法学与领域视角，Major M4）。此外，第 1 章承诺“不做文献计量统计”，但第 21 章又报告了来源池 528 条、期刊 178 条、arXiv 350 条这类计量信息；两处口径需要统一表述。

\textbf{逐节印象}。第 2 章（原理）、第 11 章（环节整合）、第 14 章（污染审计）质量最高，论证密度大；第 9 章（图神经网络）、第 8 章（生成式）偏薄，各约一页，与其在双轴矩阵中被赋予的“核心”地位不相称。

\textbf{四维小结}（该视角）：贡献—判据统一与复现审计的取向明确，是有价值的差异化，但证据呈现削弱了它；清晰—主线清楚，术语一致，个别处（核验等级）前后不一；实验强度—以合成数据演示评测纪律，方法上成立但证据力有限；完整性—六族覆盖，但深度不均。

\textbf{本视角的两条优点}：（1）判据统一，避免沦为方法清单；（2）引入“最低可接受评测协议”，有领域基础设施意义。

\textbf{本视角的核心批评}：证据表与比较矩阵的呈现方式，不足以支撑稿件自称的“可落地复现”定位。建议 major revision。

\subsection{视角 B · 方法学审稿人（研究设计与可复现，重点）}

本视角用统一清单审理，重点核查第 16 章的可复现协议与其参考实现（\texttt{code/}）是否自洽、是否可执行。

\textbf{（B-1）协议要求六类指标，实现缺一类。}第 16 章 P0 明确要求“强制报告六类指标：收益、风险、信号（IC/IR）、显著性、成本、稳健性”。但参考实现 \texttt{eval.py} 的 \texttt{metrics()} 只产出收益、风险、显著性相关量、回撤与高阶矩，\texttt{run.py} 亦未计算任何 IC 或 IR。\textbf{错在哪}：协议的“信号类”指标在实现中缺失，而 IC/IR 是截面策略最基础的诊断量。\textbf{怎么改}：在 \texttt{eval.py} 增加按期的 IC（信号与下期收益的截面相关）与 IR（IC 均值/IC 标准差），并在 \texttt{run.py} 的输出中列示；或在协议中把 IC/IR 明确为“可选”，避免协议与实现两张皮。\textbf{是否必须}：必须。

\textbf{（B-2）协议要求 walk-forward，实现未使用。}\texttt{eval.py} 提供了 \texttt{walk\_forward()} 切分函数，但 \texttt{run.py} 全程只做了一次性时间切分（前 75\% 训练、后 25\% 留出）。\textbf{错在哪}：稿件第 14 章刚批评过“单次切分容易被反复用于选模型，从而失去样本外意义”，而自己的参考实现恰好采用了它所批评的做法。这是全篇最明显的“言行不一”。\textbf{怎么改}：在 \texttt{run.py} 中用 \texttt{walk\_forward()} 组装样本外收益序列，并同时保留单次切分作为对照，展示二者结论差异。\textbf{是否必须}：必须。

\textbf{（B-3）DSR 的方差输入口径未说明。}\texttt{run.py} 以“候选池各策略夏普的截面方差”作为 \texttt{deflated\_sharpe()} 的 \texttt{var\_sr}。而 Bailey--López de Prado（2014）的 DSR 中，$\sqrt{V}$ 指\textbf{多次试验夏普估计的方差}（通常由单次夏普的标准误与试验次数推得），二者并非同一量。\textbf{错在哪}：用截面离散度替代估计方差，会使被减去的“期望最大夏普” $\text{SR}_0$ 大小失准，DSR 的数值因此不可直接解读为“校正后显著性”。\textbf{怎么改}：改用标准误口径 $\widehat{\text{SE}}(\text{SR})\approx\sqrt{(1-\gamma_3\text{SR}+\frac{\gamma_4-1}{4}\text{SR}^2)/(T-1)}$ 构造 $\sqrt{V}$，或在正文明确声明这是“以候选池离散度近似的简化”，并给出对该近似的敏感性分析。\textbf{是否必须}：必须（至少须显式声明并做敏感性分析）。

\textbf{（B-4）PBO 的切分敏感性与样本量未评估。}CSCV 取 \texttt{n\_splits}=12，在 $T=1200$ 下每段仅 100 个观测。\textbf{错在哪}：没有对 \texttt{n\_splits}（以及候选数 $M$）做敏感性分析，PBO 的单点数值（0.799）可能随切分而变。\textbf{怎么改}：报告 \texttt{n\_splits} $\in\{8,12,16,20\}$ 下的 PBO 区间，并说明其稳定性。\textbf{是否必须}：建议（major）。

\textbf{（B-5）试验次数 $N$ 的边界未定义。}\texttt{run.py} 取 $N=\text{len(cands)}=59$。但若研究者还尝试过不同的 \texttt{n\_random}、不同的岭参数 $\lambda$、不同的特征子集，真实 $N$ 将远大于 59。\textbf{错在哪}：DSR 对 $N$ 敏感，低估 $N$ 会高估显著性；“试验账本”若不覆盖全部搜索空间，就失去意义。\textbf{怎么改}：在 \texttt{run.py} 中把超参数搜索（\texttt{n\_random}、$\lambda$）全部计入账本，或在正文明确 $N$ 的定义边界并声明保守估计。\textbf{是否必须}：必须。

\textbf{（B-6）P1/P3 的协议未被实现，但未标注实现边界。}P1 要求“$\geq 3$ 个市场”的跨市场泛化，P3 要求“LLM 前向样本验证”；参考实现只跑单一合成面板，无任何跨市场与 LLM 组件。\textbf{错在哪}：读者容易把“协议”误读为“已实现能力”。\textbf{怎么改}：在 \texttt{code/} 的说明与第 16 章中给出“已实现 / 仅协议”对照表。\textbf{是否必须}：必须。

\textbf{（B-7）成本模型的两个过度简化。}冲击模型以常数 \texttt{adv}=$10^9$ 与常数日波动尺度 $k=0.02$ 建模，未随资产与时间变化；且未区分买卖方向与时段。\textbf{错在哪}：真实冲击随波动率、流动性时变，常数化会系统性低估高波动期的成本。\textbf{怎么改}：至少允许 \texttt{adv} 与 $k$ 按资产/时段给定，或声明这是量级演示。\textbf{是否必须}：建议。

\textbf{（B-8）协议遗漏交易制度约束。}协议未纳入涨跌停、停牌、退市等过滤，正文（第 11 章执行、第 16 章协议）亦几乎未讨论。\textbf{错在哪}：这些约束直接改变可交易性与成本，尤其在 A 股等市场；遗漏会使回测偏乐观。\textbf{怎么改}：在 P0 的数据/成本层加入“不可成交”掩码（涨跌停、停牌、退市）与相应处理规则。\textbf{是否必须}：必须（对面向中文市场读者的稿件尤然）。

\textbf{本视角的两条优点}：（1）把评测纪律写成可执行协议并配参考实现，方向正确且少见；（2）主动报告了不利于自己的实跑结果（样本内最优为噪声组合、成本后转负），学术诚实度高。

\textbf{本视角的核心批评}：协议与其实现之间存在多处不一致（B-1、B-2、B-6），且个别统计量（B-3）的定义未校准。建议 major revision，并在修订版中补一份“协议条款 $\leftrightarrow$ 代码实现”逐条对照表。

\subsection{视角 C · 领域审稿人（证据表与文献覆盖，重点）}

\textbf{（C-1）核验等级在正文内部不一致。}第 1 章“检索策略与证据链”约定核验等级为三档：\emph{已核验}、\emph{部分核验}、\emph{待复核}；而第 12 章证据表头与第 21 章复核清单均只使用两档（实际统计：88 条中 84 条“已核验”、4 条“部分核验”，无“待复核”）。\textbf{错在哪}：同一稿件对同一概念给出两种口径，读者无法判断“待复核”是否被使用、哪些条目属于该级。\textbf{怎么改}：统一为一套约定；若实际未使用“待复核”，则从第 1 章删除，并说明“凡不能确证元数据者不入表”。\textbf{是否必须}：必须。

\textbf{（C-2）“所有链接指回原始出处”的声明不成立。}（i）参考文献中有 4 条无链接（Almgren 等 2005、Cartea 等 2015、López de Prado 2018、Yoon 等 2019）；（ii）至少 1 条“已核验”条目的链接指向的是开放获取副本而非原始出版版本——Grossman--Stiglitz（1980）给出的是哥伦比亚大学机构库 DOI \texttt{10.7916/d8765r99}，而该文原刊为 \emph{American Economic Review} 70(3):393--408。\textbf{错在哪}：证据表的可信度依赖“可溯源到原始出处”这一承诺，个别条目已偏离该承诺。\textbf{怎么改}：对无 DOI 的经典条目（书籍、会议、杂志）明确标注“无 DOI”并给出 ISBN/会议名；对 OA 副本，注明“OA 版本链接”或改引原始出版信息。\textbf{是否必须}：必须。

\textbf{（C-3）证据表与比较矩阵口径割裂、覆盖不一致。}证据表收录 88 条，比较矩阵仅覆盖 31 项代表性工作，且两表的列结构不同（证据表无“样本外/成本口径/复现状态”，矩阵无“来源链接”）。\textbf{错在哪}：读者无法从矩阵跳到证据表核对来源，也无法从证据表判断某条目的评测严谨度；两张表各说各话，削弱了“统一口径”的主张。\textbf{怎么改}：为两表建立共同主键（如条目编号），在证据表增加“样本外/成本/复现状态”三列（或至少标注“未报告”），并在矩阵的每一行给出证据表编号。\textbf{是否必须}：必须。

\textbf{（C-4）证据表缺少“是否同行评审”与证据权重信号。}88 条中相当比例为 arXiv 预印本，但表中未标注其同行评审状态，也未给出被引量等权重线索；“核验状态=已核验”仅表示元数据互证，\textbf{与结论是否可靠无关}，极易被误读。\textbf{怎么改}：增加“来源类型（期刊/会议/预印本）”列；在表头明确写出“已核验仅指元数据层”。\textbf{是否必须}：建议（major）。

\textbf{（C-5）证据表的“关键定位”列缺可核查的填写规则。}该列为主观概括，无 rubric，不同条目详略不一（如“注意力机制”对“供应链动量传导”）。\textbf{怎么改}：给出填写规范（如限定为“该条目在本文论证中承担的角色”，并对每一族统一起见），或改为更中性的“纳入理由”。\textbf{是否必须}：建议（minor）。

\textbf{（C-6）检索过程不可独立复现。}正文给出了数据库与关键词大类，但未提供可核对的检索式清单与纳入/排除计数，也无 PRISMA 式流程图。\textbf{怎么改}：附检索式与流程图（识别 $n$ → 去重 → 纳入/排除 → 最终 $n$）。\textbf{是否必须}：建议（major）。

\textbf{（C-7）覆盖不均与市场偏置。}六族中监督学习、NLP 较厚，图网络、生成式偏薄；全篇缺少对 A 股/中文市场的专门讨论，而目标读者中相当部分来自中文市场。\textbf{怎么改}：补充中文市场小节，或明确声明范围外并给出理由。\textbf{是否必须}：建议。

\textbf{（C-8）核心术语无出处声明。}“可交易缝隙”出现在稿件标题、摘要与第 1、2、3、13 章，是贯穿全篇的组织性概念，但全文\textbf{未对它作任何来源标注}，既未说明它是工作术语，也未给出与既有概念的对应。该提法并非金融文献的通行术语（中文学术检索无相关用法）。\textbf{错在哪}：一个承担主线角色的自造术语被当作既有概念使用，读者与后续引用者会误以为它有出处；这属于出处规范问题，与 M3 同类。\textbf{怎么改}：在首次引入处（第 1 章 §1.1）显式声明为工作术语，并给出与 mispricing、profitable trading opportunity / limits to arbitrage、market impact、alpha decay 的对齐表；自推的判据须标注“本文提出”。\textbf{是否必须}：必须。

\textbf{本视角的两条优点}：（1）证据表规模（88 条）与核验痕迹（脚本 + 中间产物）在同类综述中少见；（2）比较矩阵“原文未报告则写未报告”的填表原则值得肯定。

\textbf{本视角的核心批评}：证据表自称严谨，但其口径一致性、链接规范、与矩阵的衔接，以及核心术语的出处声明，均未达到其自身标准。建议 major revision。

\subsection{视角 D · 跨学科与实际影响}

\textbf{跨学科关联}。稿件在第 2 章“一个同构性的看法”中，把交易链条与推荐系统、广告竞价的“感知—预测—决策”结构类同起来，这个类比有解释力，但未展开：三者在“环境是否对抗”“样本是否可无限采样”上的差异恰恰决定了方法迁移的边界。建议补一段形式化的对照（例如把“环境对抗性”与“样本独立性”作为两个维度，把金融定位在低样本、高对抗的象限），会显著增强跨学科说服力。

\textbf{实际影响}。第 16 章与第 15 章的实用价值最高：前者给从业者一套可执行的评测纪律，后者用实测的仓库信息（star 数、许可、最近推送）回答“能不能重跑”，这是同类综述普遍缺失的。第 20 章路线图的“负面结果库”设想有公共品价值。

\textbf{可操作性的缺口}。稿件面向“落地”，但未给出任何资源门槛的量化（算力、数据成本、人力），使读者无法判断小团队是否可执行；也未区分“个人研究者可做的部分”与“需机构资源的部分”。建议补一张“资源门槛表”。

\textbf{风险提示的充分性}。对 LLM 交易、私有另类数据的风险提示充分；但对“合规”（如另类数据涉及个人信息的法律边界）仅一句带过，考虑到读者来自受监管市场，建议单列一小段。\textbf{本视角核心批评}为建议级，不构成拒稿理由。

\subsection{视角 E · Devil's Advocate（核心论点挑战）}

\textbf{挑战一：核心命题是否可证伪？}稿件反复主张“AI 只提升转化效率，不改变利润总量约束”。但“利润总量”不可观测，命题因此难以被数据否定——一个不可证伪的命题在综述中承担主线角色，是方法论上的风险。\emph{最强反驳}：把命题弱化为可检验形式（例如“在统一口径下，跨方法族的净增量随市场成熟度递减”），并给出检验设计，否则它更像一个修辞而非一个论断。\emph{评级}：MAJOR。

\textbf{挑战二：缝隙判据是否同义反复？}命题“存在正收益调仓周期 $\iff \delta_0>c$”在数学上干净，但实践中 $\delta_0$ 与 $c$ 都需估计，而 $\delta_0$ 本身正是策略要寻找的东西。若不能给出 $\delta_0$ 的独立估计路径，该判据可能退化为“能赚钱的策略就能赚钱”。\emph{最强反驳}：需给出 $\delta_0$ 的可操作代理（如“信号在事件窗内的平均实现收益”），并说明其估计误差如何影响判据。\emph{评级}：MAJOR。

\textbf{挑战三：用合成数据演示评测纪律，是否循环论证？}稿件用一个人为低信噪比的合成面板，展示“样本内最优是纯噪声、成本后转负、PBO$=0.799$”。然而\textbf{信噪比是作者设定的}：在一个 $R^2$ 极低的设定下，“IS 选择会选到噪声”几乎是被设计出来的结果。\emph{最强反驳}：这不构成对真实市场的证据，只能作为教学演示；若要用它支撑“真实研究普遍存在此问题”的判断，必须在公开数据上做一个最小实证。稿件虽在表格标题标注“示意”，但正文措辞（“教科书式演示”）有把演示当证据之嫌。\emph{评级}：MAJOR。

\textbf{挑战四：对 LLM 的否定判断与自身标准不一致。}稿件要求 LLM 交易研究“提供前向证据”，但稿件对 LLM 交易的“低证据强度”判断，同样缺乏系统的前向证据支撑。\emph{最强反驳}：以“缺乏证据”判定“证据低”，在逻辑上成立（举证责任在主张方），但稿件应显式声明这是举证责任分配，而非实证结论。\emph{评级}：MINOR。

\textbf{Devil's Advocate 总评}：四条挑战中三条为 MAJOR，但\textbf{均非 CRITICAL}——核心命题的弱化、判据的操作化、以及合成演示的定位，都是可在修订中解决的问题。因此本视角不建议直接拒稿。

\section{阶段二：编辑综合与编辑决定}

\subsection{共识与分歧}

\textbf{五视角共识}：（1）判据统一、复现审计取向，是稿件的实质贡献（主编、领域、跨学科）；（2）证据表与比较矩阵的呈现不达其自身标准（主编、领域、方法学）；（3）协议与实现存在缺口（方法学、主编）；（4）核心术语缺出处声明（领域，编辑追认）。

\textbf{分歧与仲裁}：跨学科视角认为稿件“实用性高”，而方法学视角认为“协议未落地到实现”，二者并不矛盾——仲裁结论是：\textbf{取向实用（认可），但实现未跟上（需补）}。此外，Devil's Advocate 主张“合成演示不宜作证据”，编辑采纳其主张的精神，但认为稿件的表注与措辞已基本限定在“示意”，故降级为“须显式声明演示性质并补最小实证（可选）”。

\subsection{总评}

稿件选了一个好问题、找到了一个好切口（把评测本身作为研究对象），并在原理层给出了一条有解释力的判据。它的主要短板不在思想，而在\textbf{呈现与实现的严谨度}：证据表的口径与出处规范未达标，复现协议与参考代码多处不一致，个别统计量的定义未校准。这些都是可在一次修改中解决的问题。

\begin{table}[htbp]
\centering
\caption{编辑决定}
\begin{tabular}{@{}>{\raggedright\arraybackslash}p{3.4cm}>{\raggedright\arraybackslash}p{11.0cm}@{}}
\toprule
推荐决定 & \textbf{Major Revision（重大修改）} \\
核心理由 & 贡献明确但严谨度不足；无致命缺陷（Devil's Advocate 未判 CRITICAL） \\
\bottomrule
\end{tabular}
\end{table}

\subsection{Major comments（须逐条回应）}

\begin{longtable}{@{}>{\raggedright\arraybackslash}p{0.9cm}>{\raggedright\arraybackslash}p{5.3cm}>{\raggedright\arraybackslash}p{5.3cm}>{\raggedright\arraybackslash}p{1.9cm}@{}}
\toprule
ID & 问题（错在哪 / 在何处） & 建议解法 & 必须性 \\
\midrule
\endhead
M1 & 核验等级口径前后不一：第 1 章三档（含“待复核”），第 12/21 章两档 & 统一为一套；实际未用“待复核”则删除并说明“不可确证者不入表” & 必须 \\
M2 & 证据表与比较矩阵口径割裂、覆盖不一致（88 vs 31） & 建立共同主键；证据表增列样本外/成本/复现状态；矩阵每行标注证据表编号 & 必须 \\
M3 & “所有链接指回原始出处”不成立：4 条无链接；Grossman--Stiglitz 用 OA 副本 DOI & 无 DOI 者标注“无 DOI”并给 ISBN/会议名；OA 副本如实标注 & 必须 \\
M4 & 证据表缺“是否同行评审/来源类型”列，“已核验”易被误读为结论可靠 & 增“来源类型”列；表头明确“已核验仅指元数据层” & 必须 \\
M5 & 协议要求六类指标含 IC/IR，参考实现未实现 IC/IR & 在 eval.py/run.py 补 IC 与 IR，或协议中标为可选 & 必须 \\
M6 & 协议要求 walk-forward，run.py 只用单次切分（与第 14 章的批评自相矛盾） & 用 walk\_forward 组装 OOS，并保留单次切分对照 & 必须 \\
M7 & DSR 的 var\_sr 用候选池截面方差近似，未说明合理性 & 改用标准误口径，或声明为简化并做敏感性分析 & 必须 \\
M8 & 协议未纳入涨跌停/停牌/退市过滤 & 在 P0 数据/成本层加入不可成交掩码与处理规则 & 必须 \\
M9 & P1/P3 协议未被实现，但未标注“已实现/仅协议”边界 & 增“协议--实现”逐条对照表，并在正文标注 & 必须 \\
M10 & 试验次数 $N$ 边界未定义（未计入超参数搜索） & 账本覆盖全部搜索空间，或显式声明 $N$ 的定义边界 & 必须 \\
M11 & 核心命题“利润总量约束”不可观测、难证伪 & 弱化为可检验形式并给出检验设计 & 必须 \\
M12 & 缝隙判据 $\delta_0>c$ 缺 $\delta_0$ 的可操作估计路径 & 给出 $\delta_0$ 的代理指标与误差讨论 & 必须 \\
M13 & 合成演示的信噪比为人为设定，不能替代真实证据 & 显式声明演示性质；可选：补一个公开数据的最小实证 & 必须（声明）；建议（补实证） \\
M14 & 检索过程不可复现，无检索式清单与 PRISMA 式流程 & 附检索式与纳入/排除流程图 & 建议 \\
M15 & 核心术语“可交易缝隙”无出处声明，自造术语被当作既有概念使用（标题/摘要/§1.1/§2.4/§2.5/§3/§13） & 首次引入处声明为工作术语，附与既有概念（mispricing、limits to arbitrage、market impact、alpha decay）的对齐表；自推判据标注“本文提出” & 必须 \\
\bottomrule
\end{longtable}

\subsection{Minor comments（定位到节）}

\begin{enumerate}
  \item 第 1 章“与已有综述的关系”：把既有综述分为三类略粗，建议举 2--3 篇具体代表（含年份）以免空泛。
  \item 第 5 章：时序基础模型段落将 TimesFM 记为 2024，其 arXiv 首发为 2023；全篇年份口径（首发年 vs 发表年）应统一说明。
  \item 第 8、9 章：篇幅偏薄，与其在双轴矩阵中的“核心”标注不符，建议补各自的一节方法论细节。
  \item 第 11 章成本分解表：应显式说明“毛收益随持有期近似线性”的前提，并给出信号半衰期较短时的修正口径（现有一句提示可再强化）。
  \item 第 12 章证据表“关键定位”列：建议给出填写规范或改用更中性的“纳入理由”。
  \item 全篇：IC、IR、DSR、PBO 等缩写在首次出现处应给出定义或指向术语表；建议补术语表。
  \item 第 21 章复核清单：与第 1 章的核验等级口径同步（见 M1）。
  \item 第 15 章仓库核验表：建议注明核验时间戳，并说明 star 数等指标随时间变化、不可作为质量唯一判据。
\end{enumerate}

\subsection{修订路线图（优先级）}

\begin{enumerate}
  \item \textbf{优先一（严谨度硬伤）}：M1、M3、M5、M6、M10、M15——口径、出处与实现的一致性，改动小、收益大。
  \item \textbf{优先二（结构与呈现）}：M2、M4、M9、M14——重排两张表的主键与列，补对照表与流程图。
  \item \textbf{优先三（统计校准）}：M7、M8——修正 DSR 口径、补制度约束过滤。
  \item \textbf{优先四（论证）}：M11、M12、M13——弱化不可证伪命题、操作化判据、定位合成演示。
\end{enumerate}

\noindent\textbf{本轮修订状态注记}（据作者回执逐条核对，修订稿 68 页）：

\begin{itemize}
  \item \textbf{M1 已解决}：核验等级统一为两级，第 1 章删除“待复核”并同步第 21 章。
  \item \textbf{M2 已解决}：证据表拆为“元数据表 / 评测属性表”两张子表，以 E 编号为共同主键；比较矩阵每行标注 E 编号，两表口径由同一数据源生成，不再割裂。
  \item \textbf{M3 已解决}：无 DOI 者标注“无 DOI”，开放获取副本在“来源”列注明“OA 副本”，表头删去“所有链接指回原始出处”的绝对表述。
  \item \textbf{M4 已解决}：新增“来源类型”（期刊/会议/预印本/专著）列，表头明确“\textsc{已核验}仅指元数据层”。
  \item \textbf{M5 已解决}：\texttt{eval.py} 增加逐期截面 IC/IR，\texttt{run.py} 输出该指标。
  \item \textbf{M6 已解决}：\texttt{run.py} 改用滚动前向组装样本外序列，并保留单次切分作对照（实跑显示两者相差明显，0.477 对 0.868）。
  \item \textbf{M7 已解决}：DSR 增标准误口径（并与截面离散口径并列报告），PBO 给出四个切分段数下的区间 $[0.400,0.584]$。
  \item \textbf{M8 已解决}：P0 增加“交易制度掩码”条款，\texttt{data.py}/\texttt{baseline.py} 实现涨跌停、停牌、退市过滤。
  \item \textbf{M9 已解决}：第 16 章新增“协议条款 $\leftrightarrow$ 参考实现”对照表，逐条标注已实现 / 部分实现 / 仅协议。
  \item \textbf{M10 已解决}：试验账本计入岭参数网格，明确 $N=61$ 的构成（8 特征 + 50 随机组合 + 3 岭参数）。
  \item \textbf{M11 已解决}：核心命题弱化为可检验形式（“成本后净增量随市场效率递减”），并给出三市场同序检验设计。
  \item \textbf{M12 已解决}：给出 $\delta_0$ 与 $c$ 的可操作代理估计，并把判据改为置信下界形式。
  \item \textbf{M13 已解决}：已补“性质声明”，明确合成面板的信噪比为人为设定、只能验证流程；并新增第 21 章“真实数据最小实证”（通达信 vipdoc 日线 + AKShare 后复权因子，2019-01 至 2026-10，$T=1881$、$N=400$，同一套 P0--P2 流程）。结果与合成演示相反：DSR（标准误口径）$=0.9991$、PBO $=0.000$，说明信息层确有稳定信号；但换成可实施的仅多头口径后 OOS 夏普从 3.056 降到 1.271，成本再把多空夏普从 3.057 压到 1.164，恰好落在缝隙判据上。
  \item \textbf{M14 已解决}：第 1 章新增 PRISMA 式流程与计数（556 识别 → 528 去重 → 90 入表）及代表性检索式清单。
  \item \textbf{M15 已解决}：第 1 章 §1.1 声明“可交易缝隙”为工作术语并给出对齐表。
\end{itemize}

\noindent 八条 Minor 意见（综述分类举例、年份口径、扩充图网络与生成式两章、A 股制度小节、术语表、仓库表使用说明、核验等级同步、“关键定位”填写规范）均已处理。\textbf{排版瑕疵已消除}：证据表 A 原先的 11 处轻微溢出（均不超过 21pt）已定位为表组内 \texttt{\textbackslash emergencystretch} 被置零所致，恢复后全文档 Overfull 归零。

\noindent\textbf{结语}。这份稿件的价值在于它愿意把“评测”这一最不讨喜的部分摆到台面中央，并给出可执行协议；它需要的是把同样的诚实标准，回过头来施加到自己的证据表与代码上。修订稿已做到这一点，编辑倾向于接受。

\end{document}
